\documentclass[a4paper]{report}
\usepackage{setspace}

\pagestyle{plain}
\usepackage{amssymb}
\usepackage{graphicx}
\usepackage{color}
\usepackage{amsfonts}
\usepackage{latexsym}
\usepackage{amsmath}
\usepackage[toc,page]{appendix}
\setcounter{tocdepth}{1}
\usepackage{pdfpages}
\usepackage{todonotes}
\usepackage{hyperref}
\hypersetup{
    colorlinks,
    citecolor=black,
    filecolor=black,
    linkcolor=black,
    urlcolor=black
}

\usepackage[authoryear]{natbib}
\usepackage{algorithm}
\usepackage{algpseudocode}

% \usepackage{caption}
\usepackage{subcaption}
\usepackage{float}
\usepackage{lipsum}
\usepackage[a4paper, margin = 3cm, bottom = 2.5cm]{geometry}

\newtheorem{theorem}{THEOREM}
\newtheorem{lemma}[theorem]{LEMMA}
\newtheorem{corollary}[theorem]{COROLLARY}
\newtheorem{proposition}[theorem]{PROPOSITION}
\newtheorem{remark}[theorem]{REMARK}
\newtheorem{definition}[theorem]{DEFINITION}
\newtheorem{fact}[theorem]{FACT}

\newcommand{\nats}{\mbox{\( \mathbb N \)}}
\newcommand{\rat}{\mbox{\(\mathbb Q\)}}
\newcommand{\rats}{\mbox{\(\mathbb Q\)}}
\newcommand{\reals}{\mbox{\(\mathbb R\)}}
\newcommand{\ints}{\mbox{\(\mathbb Z\)}}

%%%%%%%%%%%%%%%%%%%%%%%%%%

% ========================================
% Title Page
% ========================================
\title{{\vspace{-14em} \includegraphics[scale=0.4]{Logos/ucl_logo.png}}\\
{{\vspace{2em} \Huge Beyond the Circle: Deforming Contours in Inverse Z-Transform}}\\
{\large Final Year Project Report}\\
}
\date{Submission date: \today}
\author{Roman Ryan Karim\thanks{
{\bf Disclaimer:}
This report is submitted as part requirement for the MEng degree in Mathematical Computation at UCL. It is substantially the result of my own work except where explicitly indicated in the text. The report may be freely copied and distributed provided the source is explicitly acknowledged.}
\\ \\ Dr Carolyn Phelan
\\ \\ \\ \\ Department of Computer Science
\\ University College London
\\ \\
}


% ========================================
% Report
% ========================================
\begin{document}
 
\onehalfspacing
\maketitle
\begin{abstract}

\end{abstract}

% ========================================
% Contents
% ========================================
\tableofcontents
\setcounter{page}{1}

% ========================================
% Introduction
% ========================================
\chapter{Introduction}
\section{Motivation}

\todo[inline]{High-Level overview for pricing options.}

\todo[inline]{How this problem relates to the inverse Z-transform.}

\section{Aims and Objectives}

\todo[inline]{Implementing widely used numerical methods for the inverse Z-transform.}
\todo[inline]{Testing out Levendorskii's method}
\todo[inline]{Benchmarking the performance of these methods.}

\section{Overview}

\todo[inline]{2-3 sentences of each chapter}

% ========================================
% Background
% ========================================
\chapter{Mathematical Background}
\section{Complex Numbers and Contours}
% Other names:
% Complex Analysis Primer
% Complex Analysis Introduction
% Intro to Complex Analysis
% Complex Analysis Basics
% Complex Analysis Fundamentals

"By definition, a complex number $z$ is an ordered pair ($x, y$) of real numbers $x$ and $y$, written $z = (x, y)$" \cite{kreyszig2010advanced}. In practice, complex numbers are written in the form $z = x + iy$, where $x$ and $y$ are real numbers and $i$ is the imaginary unit. The set of complex numbers is denoted by $\mathbb{C}$.

\section{Transform Methods}

\section{Inverse Z-Transform}

\section{Discrete Pricing Options}

% ========================================
% Experiment
% ========================================
\chapter{Experiment}

\todo[inline]{Finding different parameters to use for the experiment making use of Machine Learning techniques.}

% ========================================
% Results
% ========================================
\chapter{Results}

% ========================================
% Conclusion
% ========================================
\chapter{Conclusion}
\section{Summary}

\section{Future Work}

\section{Acknowledgements}


% ========================================
% References 
% ========================================
\addcontentsline{toc}{chapter}{References}
\bibliography{references}
\bibliographystyle{apalike}

% ========================================
% Appendix
% ========================================
\begin{appendices}

\chapter{Initial Project Plan}

\includepdf[pages=-]{initial_project_plan.pdf}
    
\end{appendices}

\end{document}